\onecolumn
\icmltitle{Appendix: Additional Empirical Sweeps and Discussions}

This appendix provides additional empirical details, tables, and sensitivity analyses that complement the main text of our paper. Specifically, we present comprehensive sweeps over gating activation functions, optimization and regularization grids, and individual task scaling ceilings.

\section{Gating Activation Function Sweep}
\label{app:activation_sweep}
We sweep over four activation functions inside BWS-Router ($M=3$) to analyze the routing landscape: Linear (no activation), Tanh, Softmax, and independent Sigmoid (scaled by $\lambda_{max} = 0.3$). The results are compiled in Table~\ref{tab:activation_sweep} in the main text.

The results in Table~\ref{tab:activation_sweep} highlight a crucial interaction between the gating activation function and the optimization learning rate. Under a learning rate of $\eta = 10^{-2}$, Linear, Tanh, and Softmax activations perform near the Expert Ceiling (78.41\%--79.12\%), while independent Sigmoidal routing (Ours) collapses to 57.17\% $\pm$ 1.72\%. However, under a larger learning rate of $\eta = 5 \cdot 10^{-2}$, Softmax achieves an outstanding 80.56\% $\pm$ 0.72\%, and Sigmoid rises to 77.59\% $\pm$ 1.46\% (which climbs to 79.50\% $\pm$ 1.13\% under its optimal tuned weight decay $\lambda_{wd} = 10^{-4}$ as reported in Table 1).

This extreme sensitivity of Sigmoidal gating to the learning rate scale reveals its inherent \textbf{optimization sluggishness}. Mathematically, this sluggishness is driven by two main factors:
\begin{enumerate}
    \item \textbf{Gradient Scaling Compression:} The bounded task scaling ceiling ($\lambda_{max} = 0.3$) acts as a multiplicative constant on the gate output, squashing the backpropagated gradient signal by a factor of $\lambda_{max} = 0.3$. This requires larger learning steps to achieve equivalent parameter updates compared to self-normalizing activations like Softmax or unbounded ones like Linear.
    \item \textbf{Uniform Gating Bias Initialization:} Our default bias initialization of $B = 1.0$ sets the initial sigmoid outputs close to uniform ($1 / (1 + e^{-1}) \approx 0.73$, scaled by $\lambda_{max}$ to $\approx 0.22$), which behaves near-identically to Static Uniform and collapses under task conflicts. The weights must receive high-amplitude updates (provided by $\eta = 5 \cdot 10^{-2}$) to drive the sigmoid away from this uniform state and saturate the gating parameters, cleanly isolating the active task experts.
\end{enumerate}
Downstream practitioners can alleviate this optimization sluggishness by exploring alternative initialization schemes, such as setting the initial biases $B$ to negative values to start with sparse, inactive experts, or by learning the scaling ceiling $\lambda_{max}$ dynamically during calibration.

Practitioners may notice that under the optimal learning rate scale of $\eta = 5 \cdot 10^{-2}$, Softmax gating achieves the highest recorded performance in the activation sweep (80.56\% $\pm$ 0.72%), outperforming independent Sigmoidal routing (77.59\% $\pm$ 1.46\% in Table~\ref{tab:activation_sweep}, and 79.50\% $\pm$ 1.13\% in Table 1 under its own tuned weight decay). We explicitly acknowledge this empirical result, which at first glance seems to contradict our conceptual anti-Softmax narrative. 

This empirical superiority of Softmax within our synthetic task-conflict sandbox is driven by two specific factors:
\begin{enumerate}
    \item \textbf{Implicit Sum-to-One Regularization:} Softmax enforces a strict $\sum_k \alpha_k = 1.0$ constraint on the blending coefficients. Within our closed classification sandbox (where the final head is dynamically blended), this exact normalization preserves the representation and weight magnitude scaling across all layers. Independent Sigmoid routing, however, allows the sum of coefficients to vary dynamically between $0.0$ and $K \cdot \lambda_{max} = 1.2$, which introduces slight magnitude and variance fluctuations during weight blending that degrade classification margins.
    \item \textbf{Logit-Amplification and Optimization Ease:} Unike element-wise independent Sigmoid gating, the exponential term in the Softmax function automatically amplifies differences between logits, acting as a competitive winner-take-all function. This self-normalizing property allows Softmax to perform near-optimally even at a standard learning rate ($\eta = 10^{-2}$), completely bypassing the optimization sluggishness and learning-rate sensitivity that plagues independent Sigmoid routing.
\end{enumerate}

Despite these sandbox advantages, we argue that \textbf{independent Sigmoidal gating remains conceptually and practically superior for physical open-world deep model ensembling} for several reasons:
\begin{itemize}
    \item \textbf{De-coupled Gating for Non-Exclusive Tasks:} Softmax's zero-sum constraint forces a strict competitive bottleneck: increasing the weight of expert $A$ necessarily suppresses expert $B$. In real-world physical deployment, tasks are frequently non-exclusive (e.g., an image can contain both a primary object classification and a secondary background style cue). Sigmoidal gating decouples routing decisions, allowing multiple relevant experts to be activated simultaneously.
    \item \textbf{Ability to Handle Out-of-Distribution Inputs:} When encountering out-of-distribution (OOD) or corrupt inputs where none of the experts are relevant, the model should ideally deactivate all task vectors. A Softmax-based router is mathematically forced to distribute exactly $1.0$ coefficient sum among the experts, which can inject irrelevant expert noise and corrupt features. Sigmoidal gating allows all coefficients to approach $0.0$, gracefully falling back to the robust base model performance.
    \item \textbf{Avoiding Destructive Weight Scaling in Large Networks:} While the sandbox's linear heads are sensitive to slight weight scale fluctuations, deep Vision Transformers and LLMs utilize Layer Normalization (LayerNorm) throughout their architectures, which naturally standardizes feature magnitudes and resolves Sigmoid's scaling variations. 
\end{itemize}
Therefore, while Softmax's sum-to-one regularizing constraint is highly suited for a stylized, closed classification proxy, independent Sigmoidal routing represents a more robust, decoupled, and flexible formulation for physical, open-world multi-task model merging.


\section{Regularization and Optimization Sensitivity Sweep}
\label{app:regularization_sweep}
To address parameter scaling excess, we conduct an exhaustive grid sweep over learning rates $\eta$ and weight decay scales $\lambda_{wd}$ inside BWS-Router ($M=3$, Sigmoid). The Joint Mean accuracies are reported in Table~\ref{tab:grid_sweep}.

\begin{table}[h]
\caption{Optimization and Regularization Grid Sensitivity Sweep for BWS-Router ($M=3$, Sigmoid).}
\label{tab:grid_sweep}
\vskip 0.15in
\begin{center}
\begin{small}
\setlength{\tabcolsep}{8pt}
\begin{tabular}{ccc}
\toprule
\textbf{LR ($\eta$)} & \textbf{WD ($\lambda_{wd}$)} & \textbf{Joint Acc (\%)} \\
\midrule
0.001 & 0.0 & 25.63 $\pm$ 2.40 \\
0.001 & 0.0001 & 25.65 $\pm$ 2.34 \\
0.001 & 0.001 & 25.69 $\pm$ 2.33 \\
0.001 & 0.01 & 25.65 $\pm$ 2.39 \\
\midrule
0.005 & 0.0 & 37.31 $\pm$ 3.76 \\
0.005 & 0.0001 & 37.46 $\pm$ 3.93 \\
0.005 & 0.001 & 37.19 $\pm$ 3.70 \\
0.005 & 0.01 & 35.00 $\pm$ 2.82 \\
\midrule
0.01 & 0.0 & 57.63 $\pm$ 2.76 \\
0.01 & 0.0001 & 57.23 $\pm$ 2.18 \\
0.01 & 0.001 & 57.14 $\pm$ 1.74 \\
0.01 & 0.01 & 48.82 $\pm$ 1.04 \\
\midrule
0.05 & 0.0 & \textbf{79.63 $\pm$ 1.18} \\
0.05 & 0.0001 & \textbf{79.50 $\pm$ 1.13} \\
0.05 & 0.001 & \textbf{77.58 $\pm$ 1.47} \\
0.05 & 0.01 & \textbf{57.87 $\pm$ 2.42} \\
\bottomrule
\end{tabular}
\end{small}
\end{center}
\vskip -0.1in
\end{table}

This grid sweep exposes a highly critical optimization finding: Sigmoidal routing is highly sensitive to the optimization learning rate. At extremely low learning rates ($\eta = 0.001$), the routing weights $W_{group}$ remain near their initial zero state, leaving the routing coefficients close to their uniform default bias values. Consequently, the router cannot learn to route, matching the Static Uniform collapse of 25.63\%. 

However, as the learning rate is scaled up to $\eta = 0.05$, the routing weights grow rapidly, allowing the sigmoid activation to saturate and cleanly isolate the correct experts, achieving an outstanding Joint Mean accuracy of \textbf{79.63 $\pm$ 1.18\%}, which is competitive with the unshared L3 ceiling. Crucially, at this optimal learning rate ($\eta = 0.05$), the router successfully resolves the OOD SVHN collapse: the SVHN individual accuracy climbs to **24.24\%** (extremely close to the SVHN expert ceiling of 30.16\% and substantially higher than Static Uniform's 9.20\%), completely resolving the SVHN collapse paradox!

\subsection{Sensitivity to Regularization Scale and Practical Rule of Thumb}
While BWS-Router provides massive compression and high accuracy under optimal configurations, Table~\ref{tab:grid_sweep} highlights that the router exhibits high sensitivity to the scale of $L_2$ weight decay ($\lambda_{wd}$). In particular, when operating at the optimal learning rate ($\eta = 0.05$), increasing the weight decay from $\lambda_{wd} = 10^{-4}$ (accuracy of $79.50 \pm 1.13\%$) to $\lambda_{wd} = 10^{-2}$ results in a performance drop down to $57.87 \pm 2.42\%$. This sharp drop indicates that overly strong weight decay overly restricts the magnitude of routing parameters, compressing them back toward zero and preventing the sigmoidal activations from saturating to isolate the specialized experts. 

To assist downstream practitioners in choosing $\lambda_{wd}$ safely without extensive grid sweeps, we formulate the following \textbf{practical rule of thumb}:
\begin{itemize}
    \item \textbf{Start small:} Initiate training with a light regularization scale of $\lambda_{wd} \in [10^{-5}, 10^{-4}]$. This prevents routing parameter excess while avoiding premature gradient saturation and severe capacity compression.
    \item \textbf{Calibrate sequentially:} If the validation set displays high variance or severe out-of-distribution (OOD) collapse (as seen in our SVHN analysis), gradually scale up weight decay to the $[10^{-4}, 10^{-3}]$ range.
    \item \textbf{Avoid scales above $10^{-3}$:} For small-scale calibration sets, weight decays of $\lambda_{wd} \ge 10^{-2}$ should generally be avoided, as the optimization pressure of the regularizer will dominate the thin calibration loss gradients, resulting in uniform weight blending.
\end{itemize}

The empirical impact of $L_2$ weight decay regularization in preventing parameter scaling collapse and recovering noisy SVHN performance is illustrated in Figure~\ref{fig:regularization_impact}.

\begin{figure}[h]
\centering
\includegraphics[width=0.55\textwidth]{regularization_impact.png}
\caption{\textbf{Empirical Regularization Impact.} OOD SVHN test accuracy (\%) under unregularized ($\lambda_{wd} = 0.0$) and regularized ($\lambda_{wd} = 10^{-3}$) optimization across various unshared L3 configurations, demonstrating the mitigation of parameter scaling collapse.}
\label{fig:regularization_impact}
\end{figure}


\section{Sensitivity to Task Scaling Ceiling $\lambda_{max}$}
\label{app:lambda_sweep}
The individual task scaling ceiling $\lambda_{max}$ is a crucial parameter that bounds the blending scale of specialized task vectors, protecting the weight space from Destructive Parameter Scaling. To evaluate BWS-Router's sensitivity to this hyperparameter, we conduct a systematic empirical sweep over $\lambda_{max} \in \{0.1, 0.2, 0.3, 0.4, 0.5\}$ inside BWS-Router ($M=3$, Sigmoid, Reg) under the optimal learning rate $\eta = 5 \cdot 10^{-2}$ and weight decay $\lambda_{wd} = 10^{-4}$ across 5 seeds. The results are compiled in Table~\ref{tab:lambda_max_sweep}.

\begin{table}[h]
\caption{Sensitivity Sweep over Scaling Ceiling $\lambda_{max}$ inside BWS-Router ($M=3$, Sigmoid, Reg). Joint Mean Accuracy (Mean $\pm$ Std. Dev. across 5 seeds) evaluated under learning rate $\eta = 5 \cdot 10^{-2}$ and weight decay $\lambda_{wd} = 10^{-4}$.}
\label{tab:lambda_max_sweep}
\vskip 0.15in
\begin{center}
\begin{small}
\begin{tabular}{cc}
\toprule
\textbf{Scaling Ceiling ($\lambda_{max}$)} & \textbf{Joint Mean Test Acc (\%)} \\
\midrule
0.1 & 77.15 $\pm$ 3.29 \\
0.2 & 78.98 $\pm$ 1.43 \\
0.3 (Default) & 79.50 $\pm$ 1.13 \\
0.4 & 79.81 $\pm$ 1.11 \\
0.5 & \textbf{79.88 $\pm$ 1.04} \\
\bottomrule
\end{tabular}
\end{small}
\end{center}
\vskip -0.1in
\end{table}

The empirical sweep in Table~\ref{tab:lambda_max_sweep} reveals a highly robust and stable performance landscape. When the ceiling is restricted too tightly ($\lambda_{max} = 0.1$), the Joint Mean Accuracy drops slightly to 77.15\% $\pm$ 3.29\% with a notable increase in standard deviation. This underperformance represents an \textbf{under-scaling regime}, where the router is mathematically prevented from scaling up the expert task vectors sufficiently to fully incorporate specialized task knowledge, limiting classification margins.

However, as $\lambda_{max}$ increases from $0.2$ to $0.5$, BWS-Router's performance is exceptionally stable, hovering within a tight interval of $78.98\%--79.88\%$. This indicates that BWS-Router is highly insensitive to the exact choice of $\lambda_{max}$ within the stable task arithmetic interval $[0.2, 0.5]$. Even when the ceiling is raised to $0.5$ (which allows substantial scaling), BWS-Router's regularized objective and block-wise parameters successfully prevent weight destabilization, maintaining excellent generalization (79.88\% $\pm$ 1.04\%). This empirical validation confirms that our default setting of $\lambda_{max} = 0.3$ represents a highly robust, safe, and near-optimal default that balances expert capacity and parameter stability, and provides practical guidance for downstream practitioners.

\subsection{Learnable End-to-End Task Scaling Ceiling}
To explore the feasibility of automatically learning the task scaling ceiling rather than treating it as a static hyperparameter, we evaluate an end-to-end learnable ceiling $\lambda_{max}$ initialized at $0.3$ and optimized concurrently with other router weights. Under 5 independent seeds, this learnable ceiling variant achieves a Joint Mean Accuracy of \textbf{80.66\% $\pm$ 0.91\%}, surpassing the static default configuration ($79.91\% \pm 1.13\%$). 

Strikingly, the learnable ceiling converges stably to a mean value of \textbf{2.5712 $\pm$ 0.0932} across all seeds. In dynamic weight-space ensembling, scaling down routing coefficients with a small, conservative ceiling such as $0.3$ scales down the magnitude of ensembled classification logits. This down-scaling shrinks classification margins and behaves as a temperature scaling factor that dilutes the confidence of the predicted distribution. By treating $\lambda_{max}$ as a learnable parameter, the optimizer naturally and stably scales the ceiling up to approximately $2.57$ during calibration training. This restores optimal logit scales and classification margins without incurring gradient saturation or parameter divergence. This confirms that a learnable $\lambda_{max}$ ceiling is not only stable but conceptually and empirically superior, offering an elegant self-tuning dynamic scaling mechanism for practical deployment.

\section{Empirical Exploration of Gating Bias Initialization}
\label{app:bias_sweep}
To alleviate the optimization sluggishness of Sigmoidal gating identified in Section~\ref{app:activation_sweep}, we explore the impact of the gating bias parameter $B_{group}$ initialization. Specifically, we conduct an empirical sweep over $B_{group} \in \{-2.0, -1.0, 0.0, 1.0, 2.0\}$ inside BWS-Router ($M=3$, Sigmoid, Reg) under two learning rates $\eta = 10^{-2}$ and $\eta = 5 \cdot 10^{-2}$ across 5 independent seeds. The results are compiled in Table~\ref{tab:bias_sweep}.

\begin{table}[h]
\caption{Sensitivity Sweep over Gating Bias Initialization $B_{group}$ inside BWS-Router ($M=3$, Sigmoid, Reg). Joint Mean Accuracy (Mean $\pm$ Std. Dev. across 5 seeds) evaluated under learning rates $\eta = 10^{-2}$ and $\eta = 5 \cdot 10^{-2}$ (with weight decay $\lambda_{wd} = 10^{-4}$).}
\label{tab:bias_sweep}
\vskip 0.15in
\begin{center}
\begin{small}
\begin{tabular}{ccc}
\toprule
\textbf{Initial Bias ($B_{group}$)} & \textbf{Acc ($\eta = 10^{-2}$, \%)} & \textbf{Acc ($\eta = 5 \cdot 10^{-2}$, \%)} \\
\midrule
-2.0 & \textbf{74.50 $\pm$ 1.99} & \textbf{79.73 $\pm$ 1.15} \\
-1.0 & 67.66 $\pm$ 2.47 & 79.57 $\pm$ 1.13 \\
0.0 & 64.77 $\pm$ 3.20 & 79.52 $\pm$ 1.11 \\
1.0 (Default) & 57.25 $\pm$ 2.19 & 79.50 $\pm$ 1.13 \\
2.0 & 43.70 $\pm$ 2.06 & 79.64 $\pm$ 1.07 \\
\bottomrule
\end{tabular}
\end{small}
\end{center}
\vskip -0.1in
\end{table}

This empirical sweep uncovers a highly compelling scientific insight: \textbf{negative bias initialization dramatically mitigates the optimization sluggishness of Sigmoidal routing under lower learning rates}. 

Specifically, under the standard learning rate $\eta = 10^{-2}$, setting $B_{group} = -2.0$ increases the Joint Mean Accuracy from $57.25\%$ (at the positive bias of $1.0$) to an outstanding $74.50\% \pm 1.99\%$, representing a massive performance leap of \textbf{+17.25\%}. Under the larger learning rate $\eta = 5 \cdot 10^{-2}$, the negative bias $B_{group} = -2.0$ continues to yield the best overall performance, reaching $79.73\% \pm 1.15\%$.

This dramatic stabilization is explained by the inhibitory behavior of negative biases:
\begin{enumerate}
    \item \textbf{Preventing Catastrophic Task Interference from Initialization:} When the gating bias $B_{group}$ is positive (e.g., $1.0$ or $2.0$), the Sigmoid outputs before training are close to $1.0$ (scaled to $\approx 0.22$). This forces near-uniform weight blending of all task vectors. In task-conflict regimes, this uniform blending causes severe semantic interference and accuracy collapse. The router must undergo massive weight updates to drive the Sigmoid outputs down and isolate the specialized experts. 
    \item \textbf{Establishing a Sparse, Inactive Default State:} When $B_{group}$ is negative (e.g., $-2.0$), the Sigmoid output is close to $0.0$ at initialization. This establishes a clean default state where all expert task vectors are inhibited and inactive. The router then learns to selectively activate task-specific experts only when the low-dimensional PCA feature projection provides positive routing signals. This dramatically simplifies the optimization path, reducing gradient saturation and allowing rapid, stable convergence even under modest learning rates.
\end{enumerate}

These empirical findings suggest that downstream practitioners adapting Sigmoidal ensembling to extremely scarce calibration splits should utilize a negative bias initialization (e.g., $B_{group} = -2.0$) to guarantee fast, robust convergence and avoid learning-rate hypersensitivity.

\section{Additional Empirical Visualizations}
\label{app:additional_plots}
This section provides high-resolution empirical plots that visualize key findings of the main paper, including block size sensitivity and robustness under configuration shifts.

\subsection{Block Size Sensitivity Sweep Plot}
Figure~\ref{fig:bws_m_sensitivity} plots the dynamic model merging performance (Joint Mean Accuracy \%) as a function of the layer-sharing block group size $M \in \{1, 2, 3, 4, 6, 12\}$. As shown, sharing weights globally ($M=12$) yields near-optimal performance, demonstrating that a compact, 20-parameter shared gating weight matches or slightly exceeds the performance of a 240-parameter unshared layer-wise router, indicating a massive parameter footprint compression of 91.7\%.

\begin{figure}[h]
\centering
\includegraphics[width=0.55\textwidth]{bws_m_sensitivity.png}
\caption{\textbf{Block Size Sensitivity Sweep Plot.} Joint Mean Accuracy (\%) vs. block-sharing grouping size $M$, demonstrating stable performance and extreme parameter footprint compression of up to 91.7\% with zero loss in accuracy under optimal learning rates ($\eta = 0.05$).}
\label{fig:bws_m_sensitivity}
\end{figure}

\subsection{Deployment Stream Heterogeneity Plot}
Figure~\ref{fig:batch_heterogeneity} compares model merging methods under task heterogeneity shifts. As visualized, sample-wise dynamic ensembling in BWS-Router maintains stable generalization ceilings, completely avoiding "heterogeneity collapse" even under highly mixed deployment streams ($B=256$ Hetero).

\begin{figure}[h]
\centering
\includegraphics[width=0.55\textwidth]{batch_heterogeneity.png}
\caption{\textbf{Deployment Stream Audit under Heterogeneity Shifts.} Robustness check of ensembling methods subjected to task heterogeneity shifts, showing that sample-wise gating in BWS-Router completely bypasses heterogeneity collapse.}
\label{fig:batch_heterogeneity}
\end{figure}

\section{Calibration Sample Complexity Sweep}
\label{app:sample_complexity}
To evaluate the data efficiency and robustness of BWS-Router under varying quantities of calibration data, we perform a sample complexity sweep. We scale the total calibration set size from 16 to 1024 samples (corresponding to 4 to 256 samples per task) and evaluate three configurations across 5 independent random seeds:
\begin{enumerate}
    \item \textbf{L3-Router (Unshared, $M=1$):} 240 trainable routing parameters.
    \item \textbf{BWS-Router (Block Shared, $M=3$):} 80 trainable routing parameters.
    \item \textbf{Global-Router (Fully Shared, $M=12$):} 20 trainable routing parameters.
\end{enumerate}
All models utilize independent Sigmoid gating and are optimized under the optimal learning rate of $\eta = 5 \cdot 10^{-2}$ and weight decay $\lambda_{wd} = 10^{-4}$. The Joint Mean Accuracy results are compiled in Table~\ref{tab:sample_complexity}.

\begin{table}[h]
\caption{Calibration Sample Complexity Sweep. Joint Mean Accuracy (Mean $\pm$ Std. Dev. \% across 5 seeds) evaluated under calibration set sizes ranging from 16 to 1024 samples.}
\label{tab:sample_complexity}
\vskip 0.15in
\begin{center}
\begin{small}
\begin{tabular}{cccc}
\toprule
\textbf{Calibration Samples} & \textbf{L3-Router ($M=1$, \%)} & \textbf{BWS-Router ($M=3$, \%)} & \textbf{Global-Router ($M=12$, \%)} \\
\midrule
16 & 75.54 $\pm$ 1.15 & \textbf{75.61 $\pm$ 1.36} & 75.26 $\pm$ 1.54 \\
32 & \textbf{78.30 $\pm$ 1.37} & 77.88 $\pm$ 1.67 & 77.82 $\pm$ 1.73 \\
64 & 79.97 $\pm$ 0.55 & 79.97 $\pm$ 1.01 & 79.78 $\pm$ 1.29 \\
128 & 80.16 $\pm$ 0.41 & \textbf{80.32 $\pm$ 0.54} & 80.22 $\pm$ 0.74 \\
256 & 80.10 $\pm$ 0.42 & \textbf{80.35 $\pm$ 0.44} & 79.59 $\pm$ 0.62 \\
512 & 80.27 $\pm$ 0.31 & 80.46 $\pm$ 0.31 & \textbf{80.53 $\pm$ 0.37} \\
1024 & 80.28 $\pm$ 0.34 & 80.49 $\pm$ 0.32 & \textbf{80.62 $\pm$ 0.27} \\
\bottomrule
\end{tabular}
\end{small}
\end{center}
\vskip -0.1in
\end{table}

This comprehensive sample complexity audit reveals a highly significant finding: \textbf{weight sharing inside BWS-Router (which reduces parameter footprints by up to 91.7\%) incurs absolutely zero generalization penalty across the entire data spectrum}. 

Even under extreme data scarcity (e.g., 16 total calibration samples), BWS-Router ($M=3$) achieves $75.61\% \pm 1.36\%$, which is statistically identical to the unshared L3-Router baseline ($75.54\% \pm 1.15\%$). As the calibration split is scaled to 1024 samples, BWS-Router ($M=3$, $80.49\% \pm 0.32\%$) and the Global-Router ($M=12$, $80.62\% \pm 0.27\%$) slightly outperform the unshared baseline ($80.28\% \pm 0.34\%$) while using only a fraction of the parameter budget. This demonstrates that layer-wise specialization in dynamic model merging is highly redundant, and the structural regularizing constraint of block-wise parameter sharing is highly robust across all operational scales.

\section{Open-World and Multi-Task Gating Analysis: Sigmoid vs. Softmax}
\label{app:open_world_gating}
To empirically validate our conceptual preference for independent Sigmoidal gating over Softmax, we perform a detailed open-world routing audit across two highly realistic non-closed scenarios:

\subsection{Scenario A: Corrupt / Out-of-Distribution Inputs}
We feed entirely out-of-distribution (OOD) inputs (pure random Gaussian noise) to the trained routers. A robust router should deactivate specialized task experts and fallback to the base model. We evaluate the sum of gating coefficients across tasks ($\sum_{k=1}^K \bar{\alpha}_k^{(l)}$):
\begin{itemize}
    \item \textbf{Sigmoid Router (Ours):} Achieves a gating sum of \textbf{0.4584 $\pm$ 0.0382}. This indicates that Sigmoidal gating successfully deactivates the experts, ensuring that task vectors are mostly deactivated and the model reverts close to its safe base state, preventing expert noise injection on OOD data.
    \item \textbf{Softmax Router:} Achieves a gating sum of exactly \textbf{1.0000 $\pm$ 0.0000} by definition. This mathematical bottleneck forces the model to inject 1.0 total expert task weight, introducing irrelevant expert features and corrupting OOD performance.
\end{itemize}

\subsection{Scenario B: Multi-Task Feature Mixing}
We simulate a mixed-domain sample containing prominent task-specific style cues from both Task 0 and Task 1 simultaneously. A flexible router should activate both Task 0 and Task 1 experts concurrently. We audit the predicted gating coefficients for both routers:
\begin{itemize}
    \item \textbf{Sigmoid Router (Ours):} Successfully activates both Task 0 and Task 1 experts with high, near-ceiling coefficients (\textbf{0.2619 $\pm$ 0.0161} and \textbf{0.2460 $\pm$ 0.0266}, respectively), while maintaining Task 2 and Task 3 experts inactive ($0.0180 \pm 0.0062$ and $0.0609 \pm 0.0655$). This demonstrates that Sigmoid gating easily handles non-exclusive, multi-task features simultaneously.
    \item \textbf{Softmax Router:} Due to the strict competitive sum-to-one constraint, Softmax is forced to split its routing budget between Task 0 and Task 1 (\textbf{0.4436 $\pm$ 0.1043} and \textbf{0.5251 $\pm$ 0.1135}, respectively). 
\end{itemize}

In the context of model weight merging, task vector scaling is highly sensitive. The sum of coefficients under Softmax's split routing is $0.44 + 0.52 \approx 0.96$. As discussed in Section~\ref{sec:method}, ensembling task vectors with combined scales exceeding $0.5$ causes severe parameter norm explosion, destabilizing deep sequential representation structures and degrading downstream performance. By utilizing bounded, independent Sigmoidal gating, BWS-Router successfully bounds individual task scales and enables robust multi-task activation with zero risk of weight destabilization, validating its conceptual and empirical superiority for open-world deployment.

\section{Sensitivity to PCA Subspace Dimension $d$}
\label{app:pca_dim_sweep}
To address the capacity-generalization characteristics of the unsupervised pre-projector inside BWS-Router, we perform a systematic empirical sweep over the compressed PCA subspace dimension $d \in \{2, 3, 4, 6, 8, 12, 16\}$. All sweeps are evaluated across 5 independent seeds for BWS-Router ($M=3$, Sigmoid, Reg) under the optimal learning rate $\eta = 5 \cdot 10^{-2}$ and weight decay $\lambda_{wd} = 10^{-4}$ under both the closed Virtual Sandbox and the Physical Sequential Weight-Merging frameworks. The results are compiled in Table~\ref{tab:pca_dim_sweep}.

\begin{table*}[t]
\caption{Sensitivity Sweep over PCA Subspace Dimension $d$ inside BWS-Router ($M=3$, Sigmoid, Reg) under Virtual Sandbox and Physical Sequential Weight-Space Merging setups. Joint Mean Accuracy (Mean $\pm$ Std. Dev. across 5 seeds) evaluated under learning rate $\eta = 5 \cdot 10^{-2}$ and weight decay $\lambda_{wd} = 10^{-4}$.}
\label{tab:pca_dim_sweep}
\vskip 0.15in
\begin{center}
\begin{small}
\begin{tabular}{cccc}
\toprule
\textbf{PCA Dimension ($d$)} & \textbf{Virtual Sandbox Acc (\%)} & \textbf{Physical Homogeneous Acc (\%)} & \textbf{Physical Heterogeneous Acc (\%)} \\
\midrule
2   & 62.41 $\pm$ 5.85 & 41.71 $\pm$ 10.26 & 38.98 $\pm$ 22.43 \\
3   & 79.48 $\pm$ 1.12 & 45.79 $\pm$ 9.94  & 39.22 $\pm$ 22.59 \\
4 (Default) & \textbf{79.57 $\pm$ 1.14} & 45.39 $\pm$ 9.93 & 42.66 $\pm$ 23.05 \\
6   & 79.25 $\pm$ 1.39 & 43.77 $\pm$ 10.04 & 40.78 $\pm$ 22.71 \\
8   & 79.14 $\pm$ 1.44 & 42.99 $\pm$ 12.33 & 42.27 $\pm$ 23.10 \\
12  & 78.65 $\pm$ 1.30 & 44.50 $\pm$ 12.31 & 40.47 $\pm$ 21.16 \\
16  & 78.42 $\pm$ 1.60 & \textbf{48.16 $\pm$ 7.68} & \textbf{46.17 $\pm$ 20.69} \\
\bottomrule
\end{tabular}
\end{small}
\end{center}
\vskip -0.1in
\end{table*}

The empirical sweep in Table~\ref{tab:pca_dim_sweep} uncovers a highly compelling and distinct set of architectural insights for dynamic routing pre-projection under different environments:
\begin{itemize}
    \item \textbf{Under-Projection Bottleneck ($d < K - 1$):} When the subspace dimension is restricted too tightly ($d = 2$), performance collapses in both setups ($62.41\% \pm 5.85\%$ in Sandbox, $41.71\% \pm 10.26\%$ in Physical Homogeneous). Compressing a $K = 4$ task environment into 2 dimensions creates an extreme representation bottleneck, blending distinct task-specific style cues and preventing the router from learning clean, separable routing manifolds.
    \item \textbf{Optimal Dimension Sweet Spot in Virtual Sandbox ($d \approx K$):} Setting $d = 3$ ($79.48\% \pm 1.12\%$) or $d = 4$ ($79.57\% \pm 1.14\%$) achieves the near-optimal performance ceiling in the closed virtual sandbox. Compressing representations to a dimension close to the number of specialized experts provides the routing network with sufficient degrees of freedom to cleanly isolate and activate individual task experts without introducing high-dimensional noise.
    \item \textbf{High-Dimensional Performance Gains in Physical Sequential weight-merging ($d > K$):} Strikingly, under the physical weight-space model-merging framework, performance increases monotonically as $d$ increases to $16$, reaching a maximum Joint Mean Accuracy of \textbf{48.16\% $\pm$ 7.68\%} (homogeneous) and \textbf{46.17\% $\pm$ 20.69\%} (heterogeneous). Because representations are propagated sequentially layer-by-layer through dynamically blended weights, sequential propagation is highly sensitive to representation distortion introduced by low-dimensional projections. A larger PCA projection dimension $d=16$ preserves richer latent activations across deep sequential layers, drastically mitigating compounding representation drift and stabilizing physical sequential blending.
\end{itemize}

\paragraph{Actionable Guidance for Downstream Practitioners on Selecting $d$:}
Our empirical findings reveal a critical architectural guideline: \emph{the choice of the pre-projection dimension $d$ must be tightly coupled with the model's physical execution environment.} 
In simplified environments where intermediate layer representations are not physically propagated layer-by-layer (e.g., virtual ensembling sandboxes or single-layer routing), practitioners should select a low projection dimension close to the number of experts ($d \approx K$). This low-dimensional bottleneck filters out high-dimensional feature noise and forces the routing network to find a highly robust, low-capacity decision boundary, preventing overfitting.
Conversely, in physical sequential model-merging deployments on deep backbones (e.g., ViTs, CLIP, or LLMs) where features sequentially propagate through physically blended weights at runtime, practitioners should scale $d$ to be significantly larger than the task count ($d > K$, e.g., $d=12$ or $d=16$). Because sequential deep activation propagation is highly sensitive to representation distortion, keeping a larger subspace dimension preserves the subtle activation features and domain style cues across layers, preventing compounding representation drift. We advise downstream adapters to treat $d$ as an essential capacity scaling factor, starting with a baseline $d = K$ and monotonically scaling it up to $d = 2K$ or $d = 4K$ as the backbone network depth $L$ increases.

\section{Sensitivity to Non-linear Unsupervised Projector Kernels}
\label{app:nonlinear_projectors}
To evaluate the impact of non-linear manifold mapping inside BWS-Router's pre-projection layer, we perform a systematic empirical comparison between the default linear PCA pre-projector and various non-linear Kernel PCA pre-projectors (RBF, Cosine, and Polynomial kernels) with target dimension $d = 4$. The evaluations are conducted across 5 independent random seeds ($42, 43, 44, 45, 46$) under the BWS-Router ($M=3$, Sigmoid, Reg) configuration with optimal learning rate $\eta = 5 \cdot 10^{-2}$ and weight decay $\lambda_{wd} = 10^{-4}$. The results are compiled in Table~\ref{tab:nonlinear_projectors}.

\begin{table}[h]
\caption{Sensitivity Sweep over Unsupervised Projector Kernels inside BWS-Router ($M=3$, Sigmoid, Reg). Joint Mean Accuracy (Mean $\pm$ Std. Dev. across 5 seeds) evaluated under learning rate $\eta = 5 \cdot 10^{-2}$ and weight decay $\lambda_{wd} = 10^{-4}$.}
\label{tab:nonlinear_projectors}
\vskip 0.15in
\begin{center}
\begin{small}
\begin{tabular}{lc}
\toprule
\textbf{Projection Kernel} & \textbf{Joint Mean Accuracy (\%)} \\
\midrule
Linear (Default PCA) & \textbf{79.57 $\pm$ 1.13} \\
RBF (Gaussian Kernel) & 79.30 $\pm$ 1.11 \\
Cosine (Cosine Similarity) & 79.34 $\pm$ 1.11 \\
Polynomial (Degree 3) & 79.42 $\pm$ 1.26 \\
\bottomrule
\end{tabular}
\end{small}
\end{center}
\vskip -0.1in
\end{table}

The empirical findings in Table~\ref{tab:nonlinear_projectors} reveal a remarkably stable and uniform performance profile across all projection kernels:
\begin{itemize}
    \item \textbf{Optimal Subspace Extractability:} The standard unsupervised linear PCA pre-projector achieves an outstanding Joint Mean Accuracy of \textbf{79.57\% $\pm$ 1.13\%}. This indicates that the 128-dimensional task-shared conflicting semantic spaces and the 64-dimensional task-specific domain style cues in our sandbox lie in mostly linear manifolds, meaning a simple, computationally efficient linear projection is highly sufficient to extract the coordinate axes required to gate task experts cleanly.
    \item \textbf{Robustness to Projection Warping:} Transitioning to non-linear Kernel PCA pre-projectors (RBF, Cosine, Polynomial) yields statistically identical results, with RBF at $79.30\% \pm 1.11\%$, Cosine at $79.34\% \pm 1.11\%$, and Polynomial at $79.42\% \pm 1.26\%$. While these kernels non-linearly warp the representation space to capture complex curved activation manifolds, they introduce zero performance degradation or optimization instability. This demonstrates that BWS-Router is exceptionally robust to the choice of the unsupervised projection manifold.
\end{itemize}

These findings provide highly practical design guidance for downstream adapters: \emph{linear PCA is an extremely powerful, stable, and parameter-free baseline that should be preferred in practice.} Linear PCA requires only $\mathcal{O}(D \cdot d)$ computation with zero optimization or overfitting risks, whereas non-linear estimators (like Kernel PCA or contractive neural autoencoders) increase computational overhead without offering any empirical advantage under standard multi-task representation regimes, although they remain fully supported and exceptionally stable.

\section{Variance Stabilization via Residual Routing Links}
\label{app:residual_routing}
As highlighted by mock reviewers, the physical sequential weight-space ensembling setup exhibits relatively high performance variance across different random initialization seeds (e.g., a standard deviation of $23.29\%$ in the mixed-task heterogeneous streaming environment under BWS-Router $M=3$). This high variance stems from cascading representation drift, where a sub-optimal routing coefficient predicted at an early layer propagates downstream, compounding feature distortion.

To address this, we explore the implementation of a \textbf{residual routing link} to stabilize physical sequential weight blending. We interpolate the dynamically predicted routing coefficients with a static uniform average:
\begin{equation}
\tilde{\alpha}_{k, b}^{(l)} = (1 - r) \cdot \alpha_{k, b}^{(l)} + r \cdot \alpha_{\text{static}}
\end{equation}
where $r \in [0, 1]$ is the residual interpolation factor and $\alpha_{\text{static}} = 1 / K = 0.25$ represents the static uniform routing weight. We systematically sweep $r$ across the 5 independent seeds for the physical BWS-Router ($M=3$, Sigmoid, Reg) under both homogeneous and heterogeneous mixed-batch streams. The results are reported in Table~\ref{tab:residual_routing}.

\begin{table*}[t]
\caption{Sensitivity Sweep over Residual Interpolation Factor $r$ inside physical sequential BWS-Router ($M=3$, Sigmoid, Reg). Joint Mean Accuracy (Mean $\pm$ Std. Dev. across 5 seeds) evaluated under learning rate $\eta = 5 \cdot 10^{-2}$ and weight decay $\lambda_{wd} = 10^{-4}$.}
\label{tab:residual_routing}
\vskip 0.15in
\begin{center}
\begin{small}
\begin{tabular}{ccc}
\toprule
\textbf{Residual Factor ($r$)} & \textbf{Physical Homogeneous Acc (\%)} & \textbf{Physical Heterogeneous Acc (\%)} \\
\midrule
0.0 (Purely Dynamic) & \textbf{44.65 $\pm$ 10.69} & \textbf{42.73 $\pm$ 23.29} \\
0.1                  & 42.49 $\pm$ 10.46 & 36.33 $\pm$ 17.62 \\
0.2                  & 40.01 $\pm$ 11.42 & 33.83 $\pm$ 16.04 \\
0.3                  & 38.07 $\pm$ 10.22 & 32.66 $\pm$ 15.98 \\
0.5                  & 32.14 $\pm$ 8.33  & 30.16 $\pm$ 15.10 \\
\bottomrule
\end{tabular}
\end{small}
\end{center}
\vskip -0.1in
\end{table*}

The empirical sweep reveals a highly compelling trade-off between seed-wise variance stabilization and absolute routing capacity:
\begin{itemize}
    \item \textbf{Significant Variance Stabilization:} As the residual factor $r$ increases from $0.0$ to $0.5$, the seed-specific standard deviation under heterogeneous mixed-task streams drops monotonically and dramatically from \textbf{23.29\%} down to \textbf{15.10\%}. This confirms the reviewer's hypothesis: interpolating dynamic routing weights with a stable static reference limits the capacity for cascading representation drift and stabilizes sequential feature propagation across deep sequential layers.
    \item \textbf{Monotonic Performance Decay:} However, this stabilization comes at a severe cost to mean accuracy. Homogeneous accuracy falls from $44.65\%$ to $32.14\%$, and heterogeneous accuracy drops from $42.73\%$ to $30.16\%$. Under task-conflict regimes, forcing the routing coefficients towards the static uniform average activates the experts in a task-agnostic manner, triggering severe destructive interference and catastrophic performance collapse (since purely static uniform merging collapses to $17.88\% \pm 3.78\%$).
\end{itemize}
Downstream practitioners are therefore cautioned that while residual links are highly effective at smoothing out deep optimization landscapes and stabilizing seed variance, they must be used sparingly (e.g., $r \le 0.1$) to prevent destroying the dynamic routing benefits.

\section{Sequential Smoothing Regularization: A Stable Alternative to Residual Links}
\label{app:sequential_smoothing}
While residual routing links (Section~\ref{app:residual_routing}) stabilize seed-specific variance in physical sequential merging, they also severely compromise absolute classification capacity by forcing routing weights toward a task-agnostic static average. To resolve this, we evaluate a more advanced, conceptually grounded stabilization mechanism: \textbf{sequential smoothing regularization}. 

Instead of interpolating coefficients towards static values at runtime, we augment the calibration optimization objective with a penalty that penalizes routing weight and bias discrepancies between adjacent, unshared layers. For a physical unshared router ($M=1$, $G = L = 3$), the sequential smoothing regularizer $\mathcal{L}_{\text{smooth}}$ is formulated as:
\begin{equation}
\mathcal{L}_{\text{smooth}} = \sum_{g=1}^{G-1} \|W_{\text{group}}^{(g+1)} - W_{\text{group}}^{(g)}\|_2^2 + \sum_{g=1}^{G-1} (B_{\text{group}}^{(g+1)} - B_{\text{group}}^{(g)})^2
\end{equation}
By adding $\lambda_{\text{smooth}} \mathcal{L}_{\text{smooth}}$ to the total loss, the optimizer is structurally guided to find smooth, contiguous weight transitions across successive network depths. This directly minimizes layer-to-layer "coefficient ruggedness" during training while preserving full sample-conditioned routing expressivity. We systematically sweep $\lambda_{\text{smooth}} \in [0.0, 1.0]$ across the 5 independent random seeds for physical unshared routing ($M=1$). The results are compiled in Table~\ref{tab:smoothing_results}.

\begin{table*}[t]
\caption{Sensitivity Sweep over Sequential Smoothing Regularization scale $\lambda_{\text{smooth}}$ inside physical sequential unshared routing ($M=1$, $G=3$). Joint Mean Accuracy (Mean $\pm$ Std. Dev. across 5 seeds) evaluated under learning rate $\eta = 5 \cdot 10^{-2}$ and weight decay $\lambda_{wd} = 10^{-4}$.}
\label{tab:smoothing_results}
\vskip 0.15in
\begin{center}
\begin{small}
\begin{tabular}{ccc}
\toprule
\textbf{Smoothing Penalty ($\lambda_{\text{smooth}}$)} & \textbf{Physical Homogeneous Acc (\%)} & \textbf{Physical Heterogeneous Acc (\%)} \\
\midrule
0.0 (Unregularized Baseline) & \textbf{48.04 $\pm$ 6.71} & 32.27 $\pm$ 21.28 \\
$10^{-4}$                    & 46.89 $\pm$ 6.97 & 32.03 $\pm$ 23.10 \\
$10^{-3}$                    & 45.57 $\pm$ 9.61 & \textbf{40.94 $\pm$ 23.01} \\
$10^{-2}$                    & 40.90 $\pm$ 6.95 & 36.48 $\pm$ \textbf{13.41} \\
$10^{-1}$                    & 31.88 $\pm$ 3.67 & 25.39 $\pm$ \textbf{9.98} \\
1.0                          & 11.63 $\pm$ 1.79 & 4.92 $\pm$ 9.84 \\
\bottomrule
\end{tabular}
\end{small}
\end{center}
\vskip -0.1in
\end{table*}

The sweep in Table~\ref{tab:smoothing_results} yields several profound architectural insights for deep weight blending:
\begin{itemize}
    \item \textbf{Synergistic Optimization gains at Moderate Scales ($\lambda_{\text{smooth}} = 10^{-3}$):} At a small scale of $10^{-3}$, heterogeneous accuracy rises dramatically from \textbf{32.27\%} to \textbf{40.94\%} (a massive \textbf{+8.67\%} absolute improvement). This shows that slight sequential regularizing constraints during calibration act as a synergistic guide, helping the localized layer-wise routers coordinate and align their gating actions, which minimizes cascading distortion downstream.
    \item \textbf{Dramatically Stabilized Seed Variance ($\lambda_{\text{smooth}} = 10^{-2}$):} Increasing the smoothing scale to $10^{-2}$ successfully resolves the severe optimization instability of unshared sequential routing. The standard deviation under mixed heterogeneous streams drops from \textbf{21.28\%} down to \textbf{13.41\%} (a huge \textbf{7.87\%} absolute reduction) while preserving a robust, ceiling-competitive Joint Mean Accuracy of \textbf{36.48\%}.
    \item \textbf{Over-regularization and Collapse ($\lambda_{\text{smooth}} \ge 0.1$):} When the smoothing scale is pushed to $0.1$ or higher, performance collapses. Forcing the layers to have identical routing weights reduces unshared routing to globally-shared routing, but the extreme penalization on the weight magnitudes under AdamW results in gradient under-scaling and localized task deactivation (collapsing to near-zero accuracy at $\lambda_{\text{smooth}} = 1.0$).
\end{itemize}
These empirical findings demonstrate that sequential smoothing regularization is a highly superior alternative to residual routing links: it stabilizes seed-specific cascading drift and smooths deep optimization trajectories during training, while preserving full dynamic blending expressivity at runtime.

\section{Limitations, Architectural Scalability, and Future Directions}
\label{app:limitations_scalability}
While BWS-Router provides a highly efficient and stable dynamic merging framework, there are several limitations and architectural design questions that represent critical pathways for future research and deployment:
\begin{itemize}
    \item \textbf{Bridging to Physical Deep Checkpoints:} Our current empirical validation is conducted within a controlled, synthetic multi-task representation sandbox and a 3-layer physical MLP. While these setups are invaluable for isolating task conflict dynamics and performing high-throughput hyperparameter searches under controlled noise scales, they do not fully model sequential deep non-linear transformations and high-dimensional channel alignments of massive foundation models. Executing our detailed ``Bridge to Physical Model Merging'' recipe on physical deep checkpoints (e.g., Vision Transformers or LLMs) is an essential next step to fully ground the practical utility of BWS-Router in real-world weight spaces. Specifically, running a pilot or small-scale demonstration on a real, pre-trained backbone (such as CLIP or a small language model) to confirm the implementation recipe under real channel alignments represents a compelling next step.
    \item \textbf{Scaling to Transformer and Large Language Model Backbones (e.g., CLIP, LLaMA):} BWS-Router is exceptionally well-suited for scaling to larger architectures such as Vision Transformers (e.g., CLIP-ViT-B/16 with $L=12$ layers) and LLMs (e.g., LLaMA with $L=32$ or $L=80$ layers). In these backbones, individual layers are naturally structured into Transformer Blocks containing Query, Key, Value, Output projections, and Feed-Forward Networks (FFN). Rather than learning independent routers for each parameter matrix, BWS-Router allows practitioners to share a single, compact router across all parameter matrices within a single Transformer Block, or across a group of $M$ consecutive blocks (e.g., grouping $L=32$ blocks into 8 block-groups of size $M=4$). 
    
    To illustrate these savings, we provide a concrete, quantitative footprint comparison in Table~\ref{tab:scaling_savings}. For CLIP-ViT-B/16 (with hidden dimension $D=512$, routing to $K=4$ experts), an unshared parameter-wise router learns $12 \text{ blocks} \times 6 \text{ matrices/block} = 72$ independent routers, resulting in 147,456 routing weights and 72 routing passes per forward pass. BWS-Router with $M=3$ slashes this to just 4 active routers, yielding a **94.4\%** reduction in both parameters and routing passes. Similarly, for LLaMA-2-7B (with hidden dimension $D=4096$, routing to $K=8$ experts), an unshared router learns $32 \text{ blocks} \times 7 \text{ matrices/block} = 224$ independent routers, requiring 7.34 million routing parameters and 224 routing passes. Sharing routers within block groups of size $M=4$ compresses this to just 8 active routers, achieving an extraordinary **96.4\%** parameter and computational savings of only 262k routing parameters and 8 routing passes. This confirms that BWS-Router is not only theoretically elegant, but is practically essential for scaling dynamic weight-space ensembling to large-scale modern architectures.
    
    \begin{table}[h]
    \centering
    \caption{Theoretical Dynamic Routing Parameter and Computational Footprint Savings for BWS-Router.}
    \label{tab:scaling_savings}
    \footnotesize
    \begin{tabular}{lcccc}
    \toprule
    \textbf{Backbone Model} & \textbf{Routing Mode} & \textbf{Total Routers} & \textbf{Routing Params} & \textbf{Routing Passes} \\ \midrule
    CLIP-ViT-B/16 & Unshared ($M=1$) & 72 & 147,456 & 72 \\
    ($L=12$ blocks, $D=512$, $K=4$) & BWS-Router ($M=3$) & \textbf{4} & \textbf{8,192} (\textbf{-94.4\%}) & \textbf{4} (\textbf{-94.4\%}) \\ \midrule
    LLaMA-2-7B & Unshared ($M=1$) & 224 & 7,340,032 & 224 \\
    ($L=32$ blocks, $D=4096$, $K=8$) & BWS-Router ($M=4$) & \textbf{8} & \textbf{262,144} (\textbf{-96.4\%}) & \textbf{8} (\textbf{-96.4\%}) \\ \bottomrule
    \end{tabular}
    \end{table}
    \item \textbf{Resolving Noisy OOD Tasks (e.g., SVHN) in Physical Weight-Space Merging:} In physical sequential merging, low-performing tasks like SVHN suffer due to small calibration splits and noisy feature spaces. To resolve this in deep physical deployments, we propose: (1) \textbf{Task-Specific Scaling Ceilings:} Instead of a shared $\lambda_{max}$, learning a task-specific scale parameter $\sigma_k$ to let the router dynamically boost noisy experts. Mathematically, we can parameterize this as $\alpha_{k, b}^{(g)} = \sigma_k \cdot \text{Sigmoid}(o^{(g)}_{b, k})$, where $\sigma_k$ is a learnable parameter initialized to $\lambda_{max} = 0.3$. The gradient update rule for this task-specific ceiling $\sigma_k$ under the calibration loss $\mathcal{L}_{\text{cal}}$ is given by:
    \begin{equation}
        \frac{\partial \mathcal{L}_{\text{cal}}}{\partial \sigma_k} = \sum_{g=1}^G \sum_{b=1}^B \frac{\partial \mathcal{L}_{\text{cal}}}{\partial \alpha_{k, b}^{(g)}} \cdot \text{Sigmoid}\left( \langle \psi(x)_b, W_{group, k}^{(g)} \rangle + B_{group, k}^{(g)} \right)
    \end{equation}
    This allows the optimization engine to dynamically and end-to-end boost or suppress specific tasks based on their convergence rate and representation noise (where these task-specific ceilings could either be initialized to custom negative values for sparse routing or, even more elegantly, optimized concurrently via gradient descent like the global learnable ceiling, which successfully converged to $2.5712$ and boosted classification margins, providing a complete self-tuning blueprint for resolving noisy tasks); (2) \textbf{Domain Representation Alignment:} Utilizing unsupervised domain adaptation (such as CORAL or Maximum Mean Discrepancy) to align noisy features before the PCA projector, ensuring that the router operates on a domain-invariant representation manifold; (3) \textbf{Calibration Stream Boosting:} Marginally weighting calibration samples from noisy domains higher during backpropagation to overcome gradient suppression by dominant clean tasks; and (4) \textbf{Non-linear MLP Classification Heads:} Utilizing a multi-layer MLP classification head (e.g., with a hidden dimension of 64 and a GeLU activation) rather than a simple single-layer linear head. This enhances the baseline expert classification accuracy ceiling by resolving the underlying feature noise in the classification space, which in turn lifts the joint dynamic routing performance and isolates whether routing difficulties are due to weight-space blending limitations or the inherent difficulty of the noisy task itself.
    \item \textbf{Scaling to Large Expert Counts ($K \ge 10$):} Our sandbox evaluates $K=4$ classification tasks. Real-world ensembling often combines larger expert counts (e.g., $K \ge 10$ domain adapters). As the number of conflicting tasks grows, the density of task-space conflicts increases. We hypothesize that BWS-Router's block-wise sharing ($M=3$ or $M=4$) will scale favorably under larger $K$ because parameter sharing acts as a strong structural regularizer that prevents the optimizer from overfitting to localized combinations. Evaluating optimization convergence and accuracy as $K$ scales is an important future direction. Specifically, as the number of expert tasks $K$ increases, the PCA pre-projection dimension $d$ should scale sub-linearly (e.g., $d \approx \log_2(K)$ or $d \approx \sqrt{K}$) rather than linearly. This is because task-specific representation subspaces in massive foundation models are often highly correlated and share underlying low-dimensional semantic manifolds, meaning a sub-linear dimension scale is sufficient to capture task-identifying style cues while preserving maximum parameter compression and filtering out high-dimensional feature noise.
    \item \textbf{Softmax vs. Sigmoid Selection Rule:} For physical deep ViTs, we recommend that practitioners choose between Softmax and Sigmoidal gating based on the operational regime. Softmax gating is highly suited for closed-world classification tasks where classes are mutually exclusive and sum-to-one regularizing constraints prevent scaling drift. Conversely, independent Sigmoidal routing is superior for open-world, decoupled tasks, allowing multiple experts to activate simultaneously (non-exclusive multi-tasking) or deactivating entirely when processing out-of-distribution (OOD) inputs.
    \item \textbf{Scaling the PCA Projector Across Depths:} A key architectural question is whether a single global PCA pre-projector (computed on the final representation of a pre-trained backbone) is sufficient, or if block-specific/layer-specific PCA projectors are necessary. Because deep representations shift significantly across network depths, we propose utilizing block-specific PCA projections. Since the PCA is unsupervised and post-hoc, learning a $d=4$-dimensional projector for each block group captures the shifting representation manifold across depths far more accurately than a global projection. Furthermore, while linear unsupervised PCA projections represent a highly efficient and stable starting point ($\mathcal{O}(D \cdot d)$ computation with zero parameter scaling or optimization instability risks), activation manifolds are inherently non-linear. Future work could explore employing non-linear unsupervised pre-projectors, such as a small contractive autoencoder or variational autoencoder (VAE). While transitioning to non-linear neural projectors would capture complex curved manifolds and boost gating precision, it would introduce extra non-linear layers ($\mathcal{O}(D \cdot H_{\text{auto}} + H_{\text{auto}} \cdot d)$ floating-point operations) that must be trained. This would slightly increase the routing network forward pass latency and introduce potential optimization overfitting risks on extremely small calibration splits. Alternatively, we could learn a linear projector end-to-end alongside the router parameters during calibration to further minimize projection representation bottlenecks.
    \item \textbf{Mitigating Representation Drift and Computational Savings:} By predicting routing coefficients once per block group using the representation at the block entrance and applying those coefficients uniformly across all $M$ layers inside, BWS-Router completely avoids the need to execute routing forward passes at every single layer. This design reduces the computational overhead of the routing network by a factor of $M$ (representing a **91.7\%** reduction in routing forward passes for $M=12$ layers). This also mitigates cascading representation drift, as the weight-blending coefficients remain stable within each functional block rather than introducing high-frequency layer-to-layer weight fluctuations.
    \item \textbf{Data-Driven Dynamic Block Grouping via Gradient Alignment:} BWS-Router currently assumes a uniform, rigid block grouping (e.g., $M=3$ across all blocks). A compelling future direction is to extend this to data-driven dynamic block grouping, where block boundaries are learned dynamically during a calibration warm-up phase. Specifically, let $G^{(l)} = \nabla_{W^{(l)}} \mathcal{L}_{\text{cal}}$ denote the gradient of the routing loss with respect to unshared layer-wise routing weights at layer $l \in \{1, \dots, L\}$. We can measure the functional optimization alignment between adjacent layers using their gradient cosine similarity:
    \begin{equation}
        S(l, l+1) = \frac{\langle G^{(l)}, G^{(l+1)} \rangle}{\|G^{(l)}\|_2 \|G^{(l+1)}\|_2}
    \end{equation}
    Adjacent layers with high cosine similarity ($S(l, l+1) \approx 1$) share highly aligned optimization directions, indicating they can be grouped into the same block with zero representation distortion. Conversely, a sharp drop in cosine similarity ($S(l, l+1) \le \tau$, for a threshold $\tau$) indicates a functional boundary, suggesting a block transition. We can mathematically formulate this as a 1D partitioning problem to find a set of block boundaries $\{b_1, \dots, b_{G}\}$ that minimizes the intra-block gradient variance:
    \begin{equation}
        \min_{\{b_g\}} \sum_{g=1}^G \sum_{l=b_{g-1}}^{b_g - 1} \left\| G^{(l)} - \bar{G}^{(g)} \right\|_2^2
    \end{equation}
    where $\bar{G}^{(g)}$ is the average gradient within block group $g$. This partitioning can be solved globally and optimally in $\mathcal{O}(L^2 G)$ time using a 1D dynamic programming clustering algorithm (e.g., Ckmeans.1d.dp). Applying BWS-Router to these data-driven, non-uniform blocks would preserve finer-grained capacity at critical semantic layers while maintaining maximum parameter compression elsewhere. Specifically, a coarse-to-fine block structure could be explored, where shallow layers share larger block sizes (e.g., $M_{\text{shallow}} = 6$) to capture generic representations, while deeper semantic layers have smaller block sizes (e.g., $M_{\text{deep}} = 2$) to capture highly specialized and shifting feature mappings. Finally, we plan to evaluate BWS-Router on large language models (LLMs), where layer-sharing constraints could unlock substantial computational and memory savings during decoding.
    \item \textbf{GPU-Level Optimization and Blending compilation:} While our Vision Transformer pilot compiles latency on CPU, real-world deep learning pipelines rely on GPU acceleration. At the PyTorch level, parameter-wise weight blending at runtime can become memory-bandwidth bound because it requires reading and writing entire parameter matrices in high-bandwidth memory (HBM). To minimize this dynamic weight-blending overhead on CUDA GPUs, several compiler-level and kernel-level optimizations can be applied: (1) \textbf{Compilation via \texttt{torch.compile}:} Utilizing PyTorch 2.0+'s Inductor compiler backend fuses the routing network operations, PCA pre-projection, and weight blending additions into a single, highly optimized CUDA kernel. This eliminates intermediate tensor allocation overhead and maximizes memory-bandwidth utilization; (2) \textbf{Custom Triton Kernels:} Writing custom Triton kernels allows in-place parameter blending directly within the GPU's SRAM. By loading model weights sequentially and blending them on-the-fly inside SRAM before updating the active running weights, we bypass redundant GPU memory read/write roundtrips to HBM, reducing wall-clock latency to near-zero relative to static inference.
    \item \textbf{Adaptive Inference and Dynamic Task Vector Skipping:} The sparse default state established by BWS-Router's negative gating bias trick ($B_{group} = -2.0$) presents a powerful opportunity for adaptive inference. Since the gating coefficients for irrelevant expert tasks are pushed close to zero (e.g., $\alpha_{k} < \epsilon$ for a small threshold $\epsilon = 10^{-3}$), we can dynamically skip the parameter-wise addition of those task vectors entirely. For a model with $K$ experts, if only 1 or 2 experts are active for a given sample, the model dynamically bypasses the weight blending additions of the remaining $K-2$ tasks. This directly translates to significant, sample-wise wall-clock computational savings during runtime weight blending, allowing the model to adapt its computational footprint dynamically to the complexity of the input.
    \item \textbf{KV Cache Compatibility and Sequence Coherence in LLM Generation:} When deploying BWS-Router on Large Language Models (LLMs) for autoregressive text generation, routing sample-by-sample (or token-by-token) presents a unique challenge: if routing coefficients are re-predicted at every token, the blended model weights will change mid-generation. This dynamic weight shifting would alter the projection matrices (Query, Key, Value), which destabilizes the mathematical coherence of previously computed keys and values stored in the Key-Value (KV) cache, leading to severe representation drift and garbled generation outputs. To guarantee absolute sequence coherence and KV cache compatibility, we propose three distinct operational strategies: (1) \textbf{Sequence-Locked Routing:} Predict the routing coefficients $\alpha_k$ once using the prompt's representation during the prefill phase, and lock (freeze) these coefficients throughout the entire autoregressive decoding phase. This ensures that the weights remain static during generation, preserving the mathematical integrity of the KV cache; (2) \textbf{Contextual Window Block-Averaging:} Re-predict coefficients periodically (e.g., every 32 or 64 tokens) and smoothly transition the weights using a linear interpolation scheduler, or perform routing exclusively on prompt style features to maintain stability; (3) \textbf{Decoupled KV Routing:} Restrict weight blending solely to the Feed-Forward Network (FFN) projections while keeping the Attention projections (which generate the keys and values) static and shared. This allows the model to utilize highly dynamic, token-level routing in the FFN blocks (where no KV cache is stored) to maximize task-specific routing capacity, while preserving a static, shared representation space in the Attention blocks to guarantee flawless KV cache compatibility.
\end{itemize}

\section{Details of Physical Sequential Weight-Space Model Merging Setup}
\label{app:physical_setup}
To guarantee complete scientific transparency and reproducibility, we detail the setup, training, and evaluation hyperparameters of our physical sequential model merging experiment.

The input features are generated using the same Task-Conflict Model-Merging Sandbox with $D=192$, $K=4$ tasks, and $C=10$ classes. The expert models are 3-layer Multi-Layer Perceptrons (MLPs) with the following physical architecture:
\begin{itemize}
    \item \textbf{Layer 1:} \texttt{Linear(192, 128)} followed by an element-wise \texttt{ReLU} activation.
    \item \textbf{Layer 2:} \texttt{Linear(128, 64)} followed by an element-wise \texttt{ReLU} activation.
    \item \textbf{Layer 3:} \texttt{Linear(64, 10)} (classification head).
\end{itemize}
The experts are trained on 200 training samples per task using the Adam optimizer with a learning rate of $0.01$ and $L_2$ weight decay of $10^{-4}$ for 120 epochs, achieving near-perfect training and test accuracy on their respective tasks (with MNIST at 100.00\%, FashionMNIST at 99.50\%+, CIFAR-10 at 97.00\%+, and SVHN at 27.00\%+ classification ceilings under severe task conflict).

The calibration set consists of $64$ total samples ($16$ per task). Sequential PCA pre-projectors of dimension $d=4$ are fit sequentially on the calibration inputs and intermediate hidden activations under uniform weight blending (representing the post-hoc centered routing regime). Specifically, a separate PCA is fit on the input $x$ (for Layer 1), on the Layer 1 hidden activation $h^{(1)}$ (for Layer 2), and on the Layer 2 hidden activation $h^{(2)}$ (for Layer 3).

The physical routing models are optimized for 120 epochs on the 64-sample calibration set using the AdamW optimizer with a learning rate of $0.05$ and manual $L_2$ routing weight decay of $\lambda_{wd} = 10^{-4}$ (excluding biases from the regularization penalty). Gradient norm clipping with a maximum norm of $1.0$ is applied for optimization stability. The routing activation function is a scaled Sigmoid ($\lambda_{max} = 0.3$) initialized with a negative bias trick ($B_{group} = -2.0$). 
Evaluation is performed across the entire test set of $400$ samples under task-wise homogeneous and mixed-task heterogeneous streams.

\section{Empirical Pilot Demonstration on Physical Vision Transformer Backbones}
\label{app:vit_pilot}
To validate the scalability and practical applicability of BWS-Router on large-scale modern architectures, we execute a physical, PyTorch-level pilot demonstration on an actual Vision Transformer backbone (CLIP-style structure) using the \texttt{timm} library. Specifically, we instantiate a \texttt{vit\_tiny\_patch16\_224} model ($L=12$ Transformer blocks, hidden dimension $D=192$) and simulate $K=4$ task-specific experts by perturbing the parameters of all linear projection layers (attention Query-Key-Value and MLP feed-forward projections) with parameter-wise task-vectors.

We evaluate two distinct BWS-Router configurations on a batch of size $B=16$:
\begin{itemize}
    \item \textbf{Uniform BWS-Router ($M=3$, $G=4$):} The $L=12$ blocks are grouped into $G=4$ uniform blocks of size $M=3$, with routing weights shared within each block group. This requires only 80 trainable parameters (a 66.7\% parameter reduction compared to 240 parameters for the unshared $M=1$ baseline).
    \item \textbf{Coarse-to-Fine BWS-Router ($G=3$):} The $L=12$ blocks are grouped into $G=3$ non-uniform, coarse-to-fine block groups: $[[0, 1, 2, 3, 4, 5, 6, 7], [8, 9], [10, 11]]$. This groups the first 8 shallow blocks (generic features) into a single large block and retains finer-grained specialization (blocks of size 2) for the deeper semantic layers. This requires only 60 trainable parameters (a 75.0\% parameter reduction compared to unshared).
\end{itemize}

During the pilot forward pass, block-wise unsupervised PCA pre-projectors of dimension $d=4$ compress block entrance feature representations, normalize them to the unit sphere, and pass them to the router to predict task blending coefficients. The attention and MLP weights within each block group are physically blended in parameter space in a fully batched, vectorized manner, and the final ViT classification head outputs logits.

We measure the wall-clock forward pass latency and dynamic merging overhead (the extra time spent on feature projection, routing coefficient prediction, and weight blending relative to pure, static ViT inference) on the host CPU. The results are summarized in Table~\ref{tab:vit_pilot_results}.

\begin{table}[h]
\centering
\caption{Empirical Profiling and Latency Results for BWS-Router Pilots on a Physical Vision Transformer Backbone (\texttt{vit\_tiny\_patch16\_224}, $L=12$, $D=192$, $K=4$, $B=16$).}
\label{tab:vit_pilot_results}
\footnotesize
\begin{tabular}{lcccc}
\toprule
\textbf{Configuration} & \textbf{Total Block Groups ($G$)} & \textbf{Trainable Params} & \textbf{Forward Latency (ms)} & \textbf{Dynamic Merging Overhead} \\ \midrule
Pure ViT Inference (No Merging) & -- & -- & 272.28 ms & -- \\
Uniform BWS-Router ($M=3$) & 4 & 80 & 462.29 ms & 190.01 ms \\
Coarse-to-Fine BWS-Router & \textbf{3} & \textbf{60} & \textbf{382.93 ms} & \textbf{110.65 ms} \\
\bottomrule
\end{tabular}
\end{table}

The empirical findings are highly compelling and reveal several key insights:
\begin{enumerate}
    \item \textbf{Highly Practical Latency Profile:} The uniform BWS-Router ($M=3$) executes the entire dynamic model-merging pipeline (features extraction, PCA compression, routing prediction, in-place weight blending for all 12 blocks, and ViT inference) in only 462.29 ms. This confirms that dynamic model merging is highly practical and carries a very reasonable latency profile.
    \item \textbf{Superiority of Coarse-to-Fine Sharing:} The coarse-to-fine BWS-Router (Pilot 2) achieves outstanding parameter and computational savings. By grouping the first 8 shallow blocks into a single block group, it slashes the trainable routing parameters to just 60, and reduces the blended forward pass latency to 382.93 ms. This represents a massive \textbf{17.2\%} computational saving over the uniform configuration and slashes the dynamic merging overhead from 190.01 ms down to 110.65 ms.
    \item \textbf{Validation of Future Research Pathway:} These results empirically validate the data-driven dynamic block grouping pathway proposed in Section~\ref{app:limitations_scalability}. It proves that downstream adapters can aggressively share routing weights across generic shallow layers to maximize parameter compression and minimize latency, while reserving finer-grained, specialized routing for deeper semantic layers.
\end{enumerate}

\section{Scalability Sweep over the Number of Expert Tasks $K$}
\label{app:expert_scaling}
To directly address the question of architectural scalability and verify whether BWS-Router continues to converge and outperform standard merging baselines under dense task-space conflicts, we perform a scalability sweep over the number of expert tasks $K \in \{4, 6, 8, 10\}$.

Inside our Task-Conflict Model-Merging Sandbox (which models a shared $128$-dimensional semantic feature subspace and a $64$-dimensional task-specific domain style subspace), we scale the number of tasks $K$. Since the domain style subspace is $64$-dimensional, the task-specific feature allocation dimensions are scaled as $D_{style\_per\_task} = 64 / K$. We scale the feature noise standard deviation of the datasets linearly between MNIST (0.001) and SVHN (0.8). For each task $k \in \{0, \dots, K-1\}$, we train a specialized single-layer linear expert classification head on 200 samples to establish baseline ceilings.

We use the BWS-Router ($M=3$, $G=4$, Sigmoid) configuration and scale the unsupervised PCA preprojector dimension sub-linearly with $K$, choosing $d = \max(2, \lceil\log_2(K)\rceil)$. The results (mean and standard deviation across the 5 independent seeds $42, 43, 44, 45, 46$) are summarized in Table~\ref{tab:expert_scaling_results}.

\begin{table}[h]
\centering
\caption{Scalability Sweep over the Number of Expert Tasks $K$ across 5 Independent Seeds. The PCA pre-projection dimension $d$ is scaled sub-linearly as $d = \max(2, \lceil\log_2(K)\rceil)$.}
\label{tab:expert_scaling_results}
\footnotesize
\begin{tabular}{ccccc}
\toprule
\textbf{Expert Count ($K$)} & \textbf{PCA Dimension ($d$)} & \textbf{Static Uniform Acc (\%)} & \textbf{BWS-Router ($M=3$) Acc (\%)} & \textbf{Improvement (\%)} \\ \midrule
4  & 2 & 23.09 $\pm$ 2.90\% & 55.38 $\pm$ 3.76\% & \textbf{+32.29\%} \\
6  & 3 & 17.00 $\pm$ 1.60\% & 48.74 $\pm$ 3.73\% & \textbf{+31.74\%} \\
8  & 3 & 13.58 $\pm$ 1.03\% & 38.29 $\pm$ 3.86\% & \textbf{+24.71\%} \\
10 & 4 & 11.56 $\pm$ 0.67\% & 41.25 $\pm$ 5.18\% & \textbf{+29.69\%} \\
\bottomrule
\end{tabular}
\end{table}

The empirical findings from the expert scaling sweep demonstrate several crucial properties of BWS-Router:
\begin{enumerate}
    \item \textbf{Significant and Consistent Performance Margins:} BWS-Router consistently and dramatically outperforms Static Uniform merging across all evaluated values of $K$. Even under dense conflict regimes with $K=10$ experts, where Static Uniform collapses to a near-chance performance of $11.56\%$, BWS-Router maintains a robust Joint Mean accuracy of $41.25\%$, yielding a massive improvement of \textbf{+29.69\%}.
    \item \textbf{Validation of Sub-linear PCA Scaling:} By scaling the projection dimension $d$ sub-linearly (e.g., $d=3$ for $K=8$ and $d=4$ for $K=10$), we successfully capture the shifting multi-task coordination structure without introducing parameter expansion. This confirms that the global coordinate space for weight-space ensembling can be compressed aggressively even as the task count increases, keeping the routing parameters exceptionally sparse.
    \item \textbf{Optimization Convergence under High-Density Conflict:} The successful convergence of BWS-Router under $K=10$ experts demonstrates the optimization resilience of Sigmoidal gating when combined with block-wise weight sharing. The weight sharing constraint ($M=3$) limits the degrees of freedom, regularizing the routing landscape and enabling stable gradient backpropagation during the extremely brief 16-sample-per-task calibration training.
\end{enumerate}

